\documentclass[aspectratio=169,11pt]{beamer}

\usetheme{Madrid}
\usecolortheme{wolverine}

\usepackage{fontspec}
\usepackage{booktabs}
\usepackage{array}
\usepackage{hyperref}
\usepackage{mihajlo-elfak-logos}

\setsansfont{Latin Modern Sans}
\setmonofont{Latin Modern Mono}
\hypersetup{colorlinks=true, urlcolor=blue}
\setbeamertemplate{navigation symbols}{}

\title{Projekat II}
\subtitle[Analitika mobilnih tokova]{Analitika velikih mobilnih tokova podataka}
\author[Mihajlo Madić 2119]{Mihajlo Madić 2119}
\institute[Big Data Systems]{Big Data Systems}
\date{\small{\textit{Jun 2026, Niš}}}

\newcommand{\code}[1]{\texttt{\small #1}}

\begin{document}

\begin{frame}
  \titlepage
\end{frame}

\begin{frame}{Cilj projekta}
  Izgraditi i uporediti dva pipeline-a za stream analitiku nad velikim SUMO skupom podataka o gradskom saobraćaju:

  \begin{itemize}
    \item Apache Spark Structured Streaming
    \item Apache Flink
  \end{itemize}

  Zajednički tok obrade:
  \[
    \text{SUMO} \rightarrow \text{Kafka} \rightarrow \text{Stream analitika} \rightarrow \text{Kafka} \rightarrow \text{Consumer / mapa}
  \]
\end{frame}

\begin{frame}{Zahtevani analitički zadaci}
  \begin{itemize}
    \item Trendovi zagađenja u blizini značajnih prostornih entiteta
    \item Intenzitet saobraćaja u blizini ulica, zona ili objekata
    \item Tumbling i sliding prozori kao runtime parametri
    \item Uporedivi rezultati iz Spark i Flink implementacije
  \end{itemize}
\end{frame}

\begin{frame}{Ciljna arhitektura}
  \begin{block}{Izvori}
    SUMO FCD izlaz, SUMO emission izlaz, referentni entiteti iz OSM / Google Maps izvora
  \end{block}

  \begin{block}{Transport}
    File-to-Kafka producer i Kafka topic-i za sirove i obrađene tokove
  \end{block}

  \begin{block}{Obrada}
    Spark Structured Streaming i Flink implementacije nad ekvivalentnim ulazima
  \end{block}

  \begin{block}{Prikaz}
    Kafka consumer sa map-oriented prikazom rezultata
  \end{block}
\end{frame}

\begin{frame}{Planirane faze}
  \begin{enumerate}
    \item Bootstrap, arhitektura i setup odluke
    \item SUMO generisanje podataka i definicija šema
    \item Kafka i kontejnerska infrastruktura
    \item Implementacija producer-a
    \item Spark analitika
    \item Flink analitika
    \item Consumer i vizualizacija
    \item Benchmarking i finalna prezentacija
  \end{enumerate}
\end{frame}

\begin{frame}{Izabrani grad}
  \begin{block}{Toulouse, France}
    Grad je fiksiran za SUMO scenario, OSM referentne entitete i primere downstream analitike.
  \end{block}

  \begin{itemize}
    \item Dovoljno veliki urbani i prigradski saobraćajni tokovi
    \item Aerodrom, glavna železnička stanica, gradski centar, stadion i bolnica
    \item Prstenasti putevi i centralni koridori pogodni za merenje intenziteta saobraćaja
  \end{itemize}
\end{frame}

\begin{frame}{Referentni prostorni entiteti}
  \begin{itemize}
    \item Snapshot: \code{data/reference/toulouse\_spatial\_entities.geojson}
    \item POI/zone: airport, Matabiau, Capitole, Stadium, Purpan
    \item Koridori: Périphérique Extérieur, Périphérique Intérieur, Route de Narbonne, Boulevard de Strasbourg
    \item Geometrije entiteta: poligoni za zone, \code{Point} za Matabiau i \code{MultiLineString} za koridore
    \item Svaki entitet ima stabilan ID, OSM reference i prag za matching
  \end{itemize}
\end{frame}

\begin{frame}{Ugovori toka i topologija}
  \begin{itemize}
    \item Kafka ulazi: \code{mobility.fcd}, \code{mobility.emissions}
    \item Kafka rezultati: \code{analytics.traffic}, \code{analytics.pollution}
    \item Spark i Flink koriste iste prozore, reference i pragove za matching
    \item Topologija: Kafka KRaft, Spark master/worker, Flink JobManager/TaskManager
    \item Run modes: \code{full}, \code{spark-only}, \code{flink-only}
    \item Bez persistent Docker volume-a u početnoj infrastrukturi
  \end{itemize}
\end{frame}

\begin{frame}{Trenutno stanje}
  \begin{itemize}
    \item Grad izabran: Toulouse, France
    \item Početni spatial reference snapshot dokumentovan sa OSM geometrijama
    \item SUMO odluka: lokalna one-shot skripta, bez kontejnerizacije za sada
    \item Kafka topic-i, event schema i početna Docker topologija implementirani
    \item Ekvivalentna spatial matching semantika u Spark i Flink aplikacijama
    \item Compose režimi: \code{full}, \code{spark-only}, \code{flink-only}
    \item Checkpoint: \code{automation/phase1\_docker\_topology.sh}
  \end{itemize}
\end{frame}

\end{document}
